\title{Controlling Text-to-Image Diffusion by Orthogonal Finetuning}

\begin{abstract}
Text-to-image diffusion models at scale exhibit remarkable abilities in translating textual descriptions into photorealistic imagery. Yet, an outstanding question remains: how can we efficiently guide or adapt these models to address diverse downstream applications? In response, we propose a systematic finetuning strategy known as Orthogonal Finetuning~(OFT) that enables flexible adaptation of text-to-image diffusion models to specialized tasks. Unlike prior approaches, OFT is theoretically guaranteed to maintain hyperspherical energy, which encodes the relational structure between neurons residing on the unit hypersphere. This invariance is found to be essential for retaining the model's semantic generative competence. For enhanced stability during finetuning, we introduce Constrained Orthogonal Finetuning (COFT), which augments OFT with an explicit radius constraint over the hypersphere. Our study focuses on two major text-to-image finetuning tasks: (1) subject-driven generation—generating images of a specific subject in novel contexts based on a limited sample and text prompt, and (2) controllable generation—allowing the model to integrate supplementary control cues. Our empirical evaluation demonstrates that OFT consistently surpasses current finetuning methodologies in both output quality and training efficiency.
\end{abstract}

\section{Introduction}

Advances in text-to-image diffusion models~\cite{saharia2022photorealistic,ramesh2022hierarchical,rombach2022high} have enabled unprecedented photorealistic rendering from textual input, exhibiting strong controllability over image synthesis. Nevertheless, relying solely on textual prompts often results in ambiguities and limited precision for nuanced image control. To mitigate these limitations, this work investigates two distinct paradigms of text-to-image generation:

\begin{itemize}[leftmargin=*,nosep]

    \item \textbf{Subject-driven generation}~\cite{ruiz2023dreambooth}: With only a handful of images depicting a particular subject, the objective is to synthesize varied depictions of that subject situated in alternative settings, leveraging textual prompts as context.
    \item \textbf{Controllable generation}~\cite{zhang2023adding,mou2023t2i}: By introducing supplementary control information (\emph{e.g.}, canny edge maps, semantic segmentation masks), this task aims to produce images that adhere to both the additional control signals and provided text.
\end{itemize}

Both of these tasks ultimately concern the challenge of finetuning text-to-image diffusion models in a manner that maintains the strong generative capabilities learned during pretraining. The essential criteria for an effective finetuning approach can be outlined as follows: (1) \emph{training efficiency}—minimizing both the number of trainable parameters and the required training epochs; and (2) \emph{generalizability preservation}—retaining the pretrained model's capacity for high-fidelity and diverse image generation. Conventionally, finetuning is conducted either by incrementally updating the network weights with a small learning rate (\emph{e.g.}, \cite{ruiz2023dreambooth}), or by introducing lightweight modules with reparameterized weights into the original model (\emph{e.g.}, \cite{hulora2022,zhang2023adding}). However, although straightforward, these methods cannot assure the safeguarding of generative performance acquired during pretraining. Furthermore, there is a gap in theoretical understanding that impedes the development of optimal finetuning strategies and the identification of appropriate hyperparameters, such as the optimal number of training epochs. A central challenge is the absence of a quantitative metric for assessing the retention of pretrained generative abilities. Most current finetuning strategies assume, without strong justification, that a smaller Euclidean distance between the finetuned and pretrained models corresponds to better preservation of generative capabilities. This belief often motivates the use of very low learning rates in practice. Nevertheless, while such an assumption may occasionally hold true, we contend that Euclidean distance is insufficient as a sole metric for semantic preservation. Therefore, developing a more structurally informed measure of the difference between the finetuned and pretrained models is likely to significantly enhance both the stability and the generative performance retention during finetuning.

Building on findings that hyperspherical similarity serves as a robust proxy for capturing semantic relationships~\cite{liu2017deep,liu2018decoupled,chen2020angular}, our approach leverages hyperspherical energy~\cite{liu2018learning} to quantitatively assess the pairwise structure of neurons within a layer. Hyperspherical energy is formally described as the aggregate of hyperspherical similarities (for instance, cosine similarities) among every pair of neurons in a layer, thus reflecting the degree of dispersion or uniformity these neurons display on the surface of the unit hypersphere~\cite{liu2021learning}. We posit that optimal finetuning should result in a minimal shift in hyperspherical energy relative to the original, pretrained model. While one straightforward strategy could introduce a regularization term to stabilize hyperspherical energy throughout finetuning, this approach does not robustly ensure minimal energy deviation. To address this, we exploit a key invariance property: when an identical orthogonal transformation is applied across all neurons, the pairwise hyperspherical similarities remain unchanged by construction. This observation underpins our proposal of Orthogonal Finetuning~(OFT), which adapts large-scale text-to-image diffusion models to downstream tasks while maintaining invariant hyperspherical energy. Fundamentally, OFT seeks a common orthogonal transformation per layer that keeps neuron pairwise angles constant, effectively recalibrating the reference frame used for neurons in that layer. Recognizing that a smaller Euclidean displacement between the finetuned and pretrained weights better conserves pretrained knowledge, we also introduce a variant, Constrained Orthogonal Finetuning~(COFT), which restricts the updated model parameters to reside within a hypersphere of fixed radius centered at the pretrained configuration.

The rationale behind employing orthogonal transformation in neuron finetuning is partly informed by insights from the 2D Fourier transform, which decomposes an image into its magnitude and phase spectra. Notably, the phase spectrum, which encodes the angular relationship between input data and the basis functions, retains most of the semantic content. To illustrate, when an image’s phase spectrum is combined with an arbitrary magnitude spectrum, it is possible to recover the image's semantic features. This observation indicates that manipulating the directions of neurons is central to achieving semantic control over generated images—a key objective in both subject-driven and controllable image generation. However, granting excessive flexibility in modifying neuron orientations risks undermining the generative capabilities acquired during pretraining. To address this, we introduce the constraint of preserving the angles between every pair of neurons. This decision is grounded in the supposition that these inter-neuronal angles are fundamental to encoding the knowledge stored in neural networks. Based on this motivation, we opt to apply a layer-wise, shared orthogonal transformation to the neurons within each layer, thereby maintaining constant hyperspherical energy.

Further conceptual motivation is drawn from orthogonal over-parameterized training~\cite{liu2021orthogonal}, a methodology in which classification networks are trained from a randomly initialized starting point via orthogonal transformation. The theoretical underpinning here is that randomly initialized networks are expected to have low hyperspherical energy, and the purpose in \cite{liu2021orthogonal} is to keep this energy minimal throughout training, as lower hyperspherical energy correlates with improved generalization in classification tasks~\cite{liu2018learning,lin2020regularizing}. \cite{liu2021orthogonal} demonstrates that orthogonal transformation affords enough adaptability to successfully train classification networks with strong generalization. In contrast, our work targets the finetuning of text-to-image diffusion models, prioritizing enhanced controllability and superior generative results on downstream tasks. We underscore two distinctions between our OFT approach and the framework of \cite{liu2021orthogonal}. First, whereas \cite{liu2021orthogonal} pursues minimization of hyperspherical energy, our OFT method is designed to maintain the hyperspherical energy characteristic of the pretrained model, thus protecting the underlying semantic representations during finetuning. Reducing this energy in diffusion models could compromise the original semantic integrity. Second, OFT introduces an explicit mechanism to limit deviations from the pretrained weights, resulting in a constrained approach—unlike \cite{liu2021orthogonal}, which operates without such restrictions. Ultimately, the effectiveness of finetuning depends on appropriately balancing adaptability with preservation of prior knowledge, and we contend that OFT achieves this equilibrium. The primary contributions of our work are as follows:

\begin{itemize}[leftmargin=*,nosep]

    \item We introduce a new finetuning technique called Orthogonal Finetuning (OFT) aimed at enhancing controllability in text-to-image diffusion models. To further ensure finetuning robustness, we also present a constrained variant that restricts the model’s angular shift relative to the original pretrained weights.
    \item In contrast to conventional finetuning strategies, OFT adjusts model parameters while analytically maintaining hyperspherical energy—an attribute we demonstrate, through experiments, to be a key factor for preserving the pretrained model’s generative semantics.
    \item We employ OFT across two primary use cases: subject-driven generation and controllable generation. A broad empirical evaluation shows that our approach outperforms previous methodologies regarding image quality, training speed, and finetuning robustness. Additionally, OFT demonstrates heightened sample efficiency, achieving strong results with fewer training examples and epochs.
    \item For the task of controllable generation, we propose a novel control consistency metric to systematically assess controllability. This metric involves inferring the control signal from the output image and comparing it to the original input signal, thus providing further evidence of OFT’s enhanced controllability.
\end{itemize}

\section{Related Work}

\textbf{Text-to-image diffusion models}. Significant advancements~\cite{saharia2022photorealistic,ramesh2022hierarchical,rombach2022high,gu2022vector,nichol2022glide} have been witnessed in the field of text-to-image generation, catalyzed largely by diffusion-based generative approaches~\cite{ho2020denoising,song2021score,song2021denoising,dhariwal2021diffusion} and the maturation of vision-language joint representation learning~\cite{su2020vlbert,lu2019vilbert,radford2021learning,wang2021simvlm,li2022blip,li2023blip,alayrac2022flamingo,schuhmann2022laion}. For example, GLIDE~\cite{nichol2022glide} and Imagen~\cite{saharia2022photorealistic} operate directly in pixel space; GLIDE employs paired text-image data for joint text encoder and diffusion prior training, whereas Imagen relies on a fixed pretrained text encoder. In contrast, latent-space approaches such as Stable Diffusion~\cite{rombach2022high} and DALL-E2~\cite{ramesh2022hierarchical} have also shown strong performance; Stable Diffusion incorporates VQ-GAN~\cite{esser2021taming} for constructing its visual latent representation, while DALL-E2 leverages CLIP~\cite{radford2021learning} for building a shared text-image latent embedding. Besides diffusion models, research on generative adversarial networks~\cite{reed2016generative,zhang2017stackgan,xu2018attngan,li2019controllable} and autoregressive models~\cite{ramesh2021zero,ding2021cogview,wu2022nuwa,yuscaling} has also been prominent in this area. OFT provides a model-agnostic finetuning strategy and is compatible with any text-to-image diffusion framework.

\textbf{Subject-driven generation}. Several works aim to restrict modifications to the subject by introducing user-provided masks as extra input~\cite{avrahami2022blended,nichol2022glide}. Approaches based on inversion~\cite{choi2021ilvr,dhariwal2021diffusion,ramesh2022hierarchical,gal2022image} enable altering image context without changing the central subject, and some, such as~\cite{hertz2022prompt}, support both local and global edits with no masking requirement. However, these typically fail to maintain identity-specific features well. Methods like Pivotal Tuning~\cite{roich2022pivotal} apply generator finetuning around an inverted latent code with identity-preserving regularization; in a similar spirit,~\cite{nitzan2022mystyle} builds a personal generative prior using a subject’s set of face images. While~\cite{casanova2021instance} can generate various forms of a specific object, it may compromise essential details specific to that instance. DreamBooth~\cite{ruiz2023dreambooth}, by incorporating a custom token and a few subject images, employs both reconstruction loss and a class prior preservation loss for finetuning text-to-image diffusion models. Distinct from straightforward low-rate finetuning, OFT adopts the DreamBooth framework but enforces orthogonal updates to the model during the finetuning process.

\textbf{Controllable generation}. Image-to-image translation falls under the broader umbrella of controllable generation, and historically, this task has been tackled primarily through the use of conditional generative adversarial networks~\cite{isola2017image,zhu2017toward,wang2018high,park2019semantic,choi2018stargan}. More recent developments have applied diffusion models to image-to-image translation~\cite{saharia2022palette,voynov2022sketch,wang2022pretraining}. In particular, ControlNet~\cite{zhang2023adding} introduced a strategy for steering pretrained diffusion models by adapting them via finetuning with extra control signals, demonstrating state-of-the-art results in controllability. Concurrently, T2I-Adapter~\cite{mou2023t2i} also leveraged finetuning of pretrained diffusion models to boost controllability over generated outputs. Building upon the experimental setup from \cite{zhang2023adding,mou2023t2i}, we utilize OFT to adapt pretrained diffusion models, achieving stronger control with reduced data and fewer parameters required for finetuning. Importantly, OFT does not result in any extra computational cost at inference time.

\textbf{Model finetuning}. Adapting large-scale pretrained models to specific downstream tasks by means of finetuning has become widespread~\cite{devlin2018bert,brown2020language,he2022masked}. In this context, adaptation techniques such as those outlined in~\cite{hulora2022,houlsby2019parameter,pfeiffer2020adapterfusion} have been explored extensively in natural language processing. The closest work to OFT is LoRA~\cite{hulora2022}, which assumes the weight updates during finetuning are of low-rank form and adds them to the existing weights. In comparison, OFT implements a layer-wise shared orthogonal transformation, updating weights multiplicatively and mathematically maintaining the relative angles between neurons within the same layer, leading to enhanced model stability.

\section{Orthogonal Finetuning}

\subsection{Why Does Orthogonal Transformation Make Sense?}

We begin our discussion by exploring the rationale behind applying orthogonal transformations in the context of finetuning text-to-image diffusion models. Specifically, we break down the inquiry into two parts: (1) the motivation for finetuning neuron angles (i.e., their orientation), and (2) the justification for employing orthogonal transformation as the mechanism for modifying these angles.

Addressing the first question, we are motivated by the empirical findings reported in \cite{liu2018decoupled,chen2020angular}, which indicate that differences in angular features effectively represent the semantic disparity. Prior work, such as SphereNet~\cite{liu2017deep}, has demonstrated that enforcing normalization of all network neurons onto a unit hypersphere not only maintains comparable model capacity but may also enhance generalization. This suggests that the neuronal direction alone can encapsulate the most salient information in the data. To further substantiate the significance of neuron angles, we perform a simple experiment, depicted in Figure~\ref{fig:toy_ae}, wherein a conventional convolutional autoencoder is trained using flower image data. During training, feature maps are computed using the standard inner product—here, $z$ represents the convolutional kernel’s element, $\bm{w}$, acting on an input within a sliding window, $\bm{x}$. At test time, we generate the feature maps in three distinct manners: (a) employing the same inner product as in training, (b) utilizing only magnitude information, and (c) isolating angular information. As illustrated in Figure~\ref{fig:toy_ae}, the angular features of neurons almost entirely reconstruct the input images, whereas the magnitude carries negligible relevant information. Importantly, the cosine activation function in scenario (c) is not used during training; the model is trained solely with the inner product. These outcomes suggest that the angles, or directions, of neurons are predominantly responsible for retaining the semantic details of the inputs. Therefore, to effectively alter the semantic content of images, adjusting the directions of neurons is likely to be the most promising approach.

Regarding the second question, the most straightforward method to refine neuronal directions is by rotating or reflecting all neurons within a layer simultaneously, a process intrinsically connected to orthogonal transformations. While more sophisticated angular transformations—such as rotating individual neurons by distinct angles—may offer greater flexibility, our findings suggest that orthogonal transformations present an optimal balance between adaptability and structural regularization. Furthermore, as established in \cite{liu2021orthogonal}, orthogonal transformations possess substantial expressiveness for training neural networks. To corroborate this viewpoint, we conduct an experiment examining the regularizing influence of orthogonality. Leveraging the controllable image generation setup of ControlNet~\cite{zhang2023adding}, we present results in Figure~\ref{fig:ortho}. The standard OFT approach demonstrates consistent performance and achieves precise control at the conclusion of training (epoch 20). Conversely, when the orthogonality constraint is omitted, OFT is unable to produce plausible images or exert meaningful control. This experiment highlights the critical role orthogonality plays within OFT.

\subsection{General Framework}

The standard approach to finetuning involves employing gradient descent with a reduced learning rate to update either the whole model or specific layers. Utilizing a small learning rate implicitly restricts the extent of deviation from the original, pretrained weights, making the typical finetuning process equivalent to training while concurrently minimizing $\thickmuskip=2mu \medmuskip=2mu \| \bm{M} - \bm{M}^0\|$, where $\bm{M}$ represents the weights post-finetuning and $\bm{M}^0$ are the initial pretrained weights. While this implicit form of regularization seems intuitively justified, it might still allow excessive freedom—especially when adapting large models. To improve upon this, LoRA explicitly imposes a low-rank constraint on the weight modifications, namely $\thickmuskip=2mu \medmuskip=2mu \text{rank}(\bm{M} - \bm{M}^0)=r'$, where $r'$ is chosen to be a small value. In contrast, OFT employs a different restriction, focusing on the neuron pairwise similarities: $\thickmuskip=2mu \medmuskip=2mu \| \text{HE}(\bm{M})-\text{HE}(\bm{M}^0) \|=0$, with $\text{HE}(\cdot)$ representing the hyperspherical energy of a weight matrix. 

To exemplify, consider a fully connected layer $\thickmuskip=2mu \medmuskip=2mu \bm{W}=\{\bm{w}_1,\cdots,\bm{w}_n\}\in\mathbb{R}^{d\times n}$, where the vector $\bm{w}_i\in\mathbb{R}^{d}$ denotes the $i$-th neuron's weights (with $\bm{W}^0$ standing for the pretrained version). The resulting output vector $\thickmuskip=2mu \medmuskip=2mu \bm{z}\in\mathbb{R}^n$ is determined by $\thickmuskip=2mu \medmuskip=2mu \bm{z}=\bm{W}^\top\bm{x}$, for input $\thickmuskip=2mu \medmuskip=2mu \bm{x}\in\mathbb{R}^d$. The goal of OFT is to reduce the difference in hyperspherical energy between the updated and the original models, as given by

\begin{equation}\label{eq:oft_principle}
\footnotesize
\min_{\bm{W}} \left\| \text{HE}(\bm{W})-\text{HE}(\bm{W}^0)\right\|~~\Leftrightarrow~~\min_{\bm{W}} \bigg{\|} \sum_{i\neq j} \|\hat{\bm{w}}_i-\hat{\bm{w}}_j\|^{-1}-\sum_{i\neq j} \|\hat{\bm{w}}^0_i-\hat{\bm{w}}^0_j\|^{-1}\bigg{\|}
\end{equation}

where $\thickmuskip=2mu \medmuskip=2mu \hat{\bm{w}}_i:=\bm{w}_i/\|\bm{w}_i\|$ designates the normalized weights of neuron $i$, and hyperspherical energy for $\bm{W}$ is formulated as $\thickmuskip=2mu \medmuskip=2mu \text{HE}(\bm{W}):= \sum_{i\neq j} \|\hat{\bm{w}}_i-\hat{\bm{w}}_j\|^{-1}$. Notably, the optimal solution to Eq.~\eqref{eq:oft_principle} achieves zero, which occurs when $\bm{W}$ and $\bm{W}^0$ are connected by a rotation or reflection, i.e., $\thickmuskip=2mu \medmuskip=2mu \bm{W}=\bm{R}\bm{W}^0$, with $\thickmuskip=2mu \medmuskip=2mu \bm{R}\in\mathbb{R}^{d\times d}$ as an orthogonal matrix (where the determinant $1$ or $-1$ indicates a rotation or a reflection, respectively). This fundamental principle underpins OFT: to perform finetuning by learning an orthogonal matrix for each layer, transforming all neurons within a

\begin{equation}\label{eq:oft_general}
    \footnotesize
    \bm{z}=\bm{W}^\top\bm{x}=(\bm{R}\cdot\bm{W}^0)^\top\bm{x},~~~\text{s.t.}~\bm{R}^\top\bm{R}=\bm{R}\bm{R}^\top=\bm{I}
\end{equation}

In this formulation, $\bm{W}$ corresponds to the weight matrix obtained via OFT finetuning, and $\bm{I}$ stands for the identity matrix. Figure~\ref{fig:block} provides an overview of OFT. Analogous to the zero initialization used in LoRA, it is necessary for OFT to update the pretrained model starting exactly from $\bm{W}^0$. This is accomplished by setting the initial value of the orthogonal matrix $\bm{R}$ as the identity matrix, which ensures that the starting point for finetuning coincides with the pretrained parameters. To enforce the orthogonality condition on $\bm{R}$, differentiable orthogonalization techniques such as those described in \cite{liu2021orthogonal,lezcano2019cheap} may be leveraged. We will further address strategies to maintain orthogonality efficiently in the following discussion.

\subsection{Efficient Orthogonal Parameterization}\label{sect:eff_ortho}

While differentiable, typical orthogonalization processes like the Gram-Schmidt algorithm tend to be computationally prohibitive for practical use~\cite{liu2021orthogonal}. To improve computational efficiency, we employ Cayley parameterization as a means to generate orthogonal matrices. More specifically, we define the orthogonal matrix as $\thickmuskip=2mu \medmuskip=2mu \bm{R}=(\bm{I}+\bm{Q})(I-\bm{Q})^{-1}$, where $\bm{Q}$ is a skew-symmetric matrix, that is, $\thickmuskip=2mu \medmuskip=2mu \bm{Q}=-\bm{Q}^\top$. This approach provides significant efficiency gains but with a minor restriction—Cayley parameterization can only yield orthogonal matrices whose determinant is $1$, confining $\bm{R}$ to the special orthogonal group. In practice, however, this constraint does not negatively impact the results. Despite utilizing the Cayley transform for orthogonal parameterization, $\bm{R}$ may remain parameter-inefficient if the dimension $d$ is large. To mitigate this, we propose representing $\bm{R}$ as a block-diagonal matrix composed of $r$ smaller blocks, as given below:

\begin{equation}
    \footnotesize
    \bm{R}=\text{diag}(\bm{R}_1,\bm{R}_2,\cdots,\bm{R}_r)=\begin{bmatrix}
    \bm{R}_1\in \text{O}(\frac{d}{r}) &  &\\
    &\ddots & \\
    & & \bm{R}_r\in \text{O}(\frac{d}{r})
    \end{bmatrix}\in \text{O}(d)
\end{equation}

where $\text{O}(d)$ represents the orthogonal group in $d$-dimensional space, with $\thickmuskip=2mu \medmuskip=2mu \bm{R}\in\mathbb{R}^{d\times d}$ and $\thickmuskip=2mu \medmuskip=2mu \bm{R}_i\in\mathbb{R}^{d/r\times d/r}$ for all $i$ denoting orthogonal matrices. When $\thickmuskip=2mu \medmuskip=2mu r=1$, the block-diagonal orthogonal matrix simply reduces to a regular, unconstrained orthogonal matrix. An orthogonal matrix of dimension $\thickmuskip=2mu \medmuskip=2mu d\times d$ requires $\thickmuskip=2mu \medmuskip=2mu d(d-1)/2$ independent parameters, which corresponds to a parameter complexity of $\mathcal{O}(d^2)$. In contrast, for an $r$-block diagonal orthogonal structure, the parameter count is $\thickmuskip=2mu \medmuskip=2mu d(d/r-1)/2$, yielding a reduced complexity of $\thickmuskip=2mu \medmuskip=2mu \mathcal{O}(d^2/r)$. An additional reduction in parameter count is achievable by tying all block matrices together, that is, by setting $\thickmuskip=2mu \medmuskip=2mu\bm{R}_i=\bm{R}_j$ for every $i\neq j$, thus lowering the parameter complexity further to $\mathcal{O}(d^2/r^2)$. Importantly, these efficiency gains do not compromise the orthogonality of $\bm{R}$, ensuring that hyperspherical energy is strictly conserved.

Next, we examine the parameter usage of OFT compared to LoRA. For a LoRA configuration with rank $r'$, its number of learnable parameters is $\thickmuskip=2mu \medmuskip=2mu r'(d+n)$. If both $r$ and $r'$ scale proportionally with $d$, i.e., $\thickmuskip=2mu \medmuskip=2mu r = r' = \alpha d$ for some constant $\thickmuskip=2mu \medmuskip=2mu 0<\alpha\leq 1$, the complexity of LoRA becomes $\mathcal{O}(d^2+dn)$, whereas for OFT it scales as $\mathcal{O}(d)$. To concretely highlight the parameter savings, consider a weight matrix of size $\thickmuskip=2mu \medmuskip=2mu 128\times 128$: with $\thickmuskip=2mu \medmuskip=2mu r'=8$, LoRA requires $2,048$ parameters, while OFT needs only $960$ parameters for $\thickmuskip=2mu \medmuskip=2mu r=8$ (without block sharing).

\subsection{Constrained Orthogonal Finetuning}

The flexibility inherent in OFT can be restricted further by confining the adaptation to a small region near the original, pretrained model. In particular, COFT implements the following form for the forward computation:

\begin{equation}\label{eq:coft_general}
    \footnotesize
    \bm{z}=\bm{W}^\top\bm{x}=(\bm{R}\cdot\bm{W}^0)^\top\bm{x},~~~\text{s.t.}~\bm{R}^\top\bm{R}=\bm{R}\bm{R}^\top=\bm{I},~\norm{\bm{R}-\bm{I}}\leq \epsilon
\end{equation}

where there is not only an orthogonality requirement but also a constraint limiting the deviation from the identity matrix to within $\epsilon$. Orthogonality may be enforced using the Cayley parameterization described previously in Section~\ref{sect:eff_ortho}, but enforcing the $\epsilon$-neighborhood restriction with Cayley-parameterized orthogonal matrices is nontrivial. To better understand this constraint, we apply the Neumann series to the Cayley transform $\thickmuskip=2mu \medmuskip=2mu \bm{R}=(\bm{I}+\bm{Q})(I-\bm{Q})^{-1}$ and approximate as $\thickmuskip=2mu

Since the series converges in the operator norm, it becomes possible to impose the constraint $\thickmuskip=2mu \medmuskip=2mu\|\bm{R}-\bm{I}\|\leq\epsilon$ within the Cayley transform. The constraint is equivalently reformulated as $\thickmuskip=2mu \medmuskip=2mu \|\bm{Q}-\bm{0}\|\leq\epsilon'$, where $\bm{0}$ is a zero matrix and $\epsilon'$ denotes a separate tolerance parameter distinct from $\epsilon$. Enforcing this updated restriction on $\bm{Q}$ is straightforward using projected gradient descent. Initializing the orthogonal matrix $\bm{R}$ as an identity matrix corresponds to setting $\bm{Q}$ to the zero matrix. Conceptually, COFT explicitly unifies two types of constraints: one constraining the minimal hyperspherical energy change, and another restricting deviation from the original pretrained model. While the latter is often adopted implicitly by existing finetuning schemes, COFT formalizes it as an explicit objective. Although OFT delivers robust performance, empirical results indicate that COFT enhances finetuning stability further due to this precise control over model deviation. Figure~\ref{fig:eps_coft} illustrates the influence of $\epsilon$ on COFT performance, demonstrating that $\epsilon$ modulates finetuning adaptability. For larger values of $\epsilon$, COFT-optimized models closely resemble those obtained with OFT; for smaller $\epsilon$, the COFT-adjusted model converges towards the behavior of the original pretrained text-to-image diffusion system.

\subsection{Re-scaled Orthogonal Finetuning}

We introduce a straightforward augmentation of OFT, in which a scaling factor is learned for each neuron independently. The rationale stems from the observation that rescaling neurons has no impact on the hyperspherical energy, since magnitude effects are normalized. Concretely, we adopt the forward computation $\thickmuskip=2mu \medmuskip=2mu\bm{z}=(\bm{R}\bm{W}^0\bm{D})^\top\bm{x}$\footnote{\textbf{Errata}: In the NeurIPS camera-ready submission, the re-scaled OFT forward pass was mistakenly presented as $\thickmuskip=2mu \medmuskip=2mu\bm{z}=(\bm{D}\bm{R}\bm{W}^0)^\top\bm{x}$. The implementation used for experiments is correct, and therefore the findings in Appendix~\ref{app:rescaled} are accurate.} where $\thickmuskip=2mu \medmuskip=2mu \bm{D}=\text{diag}(s_1,\cdots,s_n)\in\mathbb{R}^{n\times n}$ is a diagonal matrix with all entries $s_1,\cdots,s_n$ learnable and strictly positive. Unlike the original OFT forward mapping from Eq.~\eqref{eq:oft_general}, which optimizes solely over $\bm{R}$, this approach allows both the diagonal scaling matrix $\bm{D}$ and the orthogonal matrix $\bm{R}$ to be learned. The resulting re-scaled OFT thus enhances the representational capacity of OFT, requiring minimal extra parameters. For the sake of evaluating the role of orthogonal transformations specifically, we retain standard OFT for the main experiments; nevertheless, additional tests (see Appendix~\ref{app:rescaled}) suggest that re-scaled OFT generally offers superior performance.

\section{Intriguing Insights and Discussions}

\textbf{OFT exhibits architectural generality}. In theory, OFT is applicable across all neural network architectures. For example, in Transformers, LoRA is generally integrated with attention mechanisms~\cite{hulora2022}. To fairly benchmark OFT against LoRA, our empirical setup confines OFT updates to the attention weights as well. Moreover, in addition to use with fully connected networks, OFT is particularly apt for convolutional layers, since the block-diagonal nature of $\bm{R}$ has interpretability advantages not present in LoRA. With block partitions matched to the count of input channels, each block specifically acts on a distinct channel, analogous to learning separate depth-wise convolutional filters~\cite{chollet2017xception}. Alternatively, sharing all blocks of $\bm{R}$ implies the same orthogonal transformation is applied across all channels. Initial explorations on adapting OFT for convolutional layers are documented in Appendix~\ref{app:convolution}.

\textbf{Connection to LoRA}. LoRA stabilizes the pretrained weight matrix by incorporating a low-rank matrix, thereby limiting significant deviations in the existing parameters. In contrast, OFT employs a transformation constrained to be orthogonal (and therefore full-rank), ensuring that the alteration to the pretrained matrix does not eliminate previously learned knowledge. The forward computation for OFT can be reformulated as $\thickmuskip=2mu \medmuskip=2mu \bm{z}=(\bm{R}\bm{W}^0)^\top\bm{x}=(\bm{W}^0+(\bm{R}-\bm{I})\bm{W}^0)^\top\bm{x}$, where $\thickmuskip=2mu \medmuskip=2mu (\bm{R}-\bm{I})\bm{W}^0$ serves a role analogous to the low-rank correction utilized in LoRA. Given that $\bm{W}^0$ is usually full-rank, OFT effectively becomes a low-rank update mechanism whenever $\thickmuskip=2mu \medmuskip=2mu \bm{R}-\bm{I}$ has low-rank structure. Paralleling LoRA's introduction of a rank parameter $r'$ to regulate the number of trainable variables, OFT controls trainability through the diagonal block size $r$. Fundamentally, LoRA and OFT each pursue parameter efficiency by distinct means: LoRA capitalizes on the inherent low-rank features to reduce the number of tunable weights, whereas OFT leverages sparsity, specifically via block-diagonal orthogonality.

\textbf{Why OFT converges faster}. Firstly, evidence from Figure~\ref{fig:toy_ae} indicates that semantic modifications are most efficiently achieved by altering neuron orientations, which aligns with the primary design of OFT. Secondly, OFT's optimization can be interpreted as fine-tuning over a smooth hypersphere manifold, resulting in a more tractable optimization landscape. Empirical support for this perspective can also be found in \cite{liu2021orthogonal}.

\textbf{Why not minimize hyperspherical energy}. Distinct from \cite{liu2021orthogonal}, our approach does not target minimization of hyperspherical energy. In the context of classification, it is advantageous for neurons to exhibit no redundancy, but enforcing minimal hyperspherical energy would cause neurons to be evenly distributed over the hypersphere. Such uniform spacing is not suitable for finetuning purposes because it can compromise the information encoded during pretraining.

\textbf{Balancing flexibility and regularity in finetuning}. Our findings highlight a fundamental tension between the flexibility and regularity of finetuning approaches. Conventional finetuning techniques offer maximum adaptability, yet this often comes at the expense of model stability and increases the likelihood of model collapse. OFT, despite its simplicity, achieves a desirable compromise by maintaining pairwise neuron angle relationships, thus mediating between the two extremes. The block-diagonal parameterization, furthermore, can be interpreted as a form of intensified regularization on the orthogonal matrix.

\textbf{Zero inference-time cost}. Contrary to approaches such as ControlNet, the OFT framework does not impose any extra computational burden at inference. Specifically, during the inference phase, the learned orthogonal transformation $\bm{R}$ can be directly composed with the original pretrained weight matrix $\bm{W}^0$, forming the equivalent weight $\thickmuskip=2mu \medmuskip=2mu \bm{W}=\bm{R}\bm{W}^0$. Consequently, the runtime performance is identical to that of the pretrained architecture.

\section{Experiments and Results}

\textbf{Experimental configuration}. All experiments utilize Stable Diffusion v1.5~\cite{rombach2022high} as the backbone text-to-image model. For unbiased evaluation, images produced by each method are randomly sampled. Subject-driven generation experiments adhere to the DreamBooth~\cite{ruiz2023dreambooth} protocol, while controlled generation is conducted according to ControlNet~\cite{zhang2023adding} and T2I-Adapter~\cite{mou2023t2i} standards. In making comparisons with LoRA, we restrict OFT or COFT modifications to the same layers affected by LoRA for fair benchmarking. Expanded experimental details and additional results appear in Appendix~\ref{app:settings}.

\subsection{Subject-driven Generation}

\textbf{Experiment details}. For baselines, we employ both DreamBooth~\cite{ruiz2023dreambooth} and LoRA~\cite{hulora2022}. Across all approaches, the loss function from DreamBooth is applied consistently. DreamBooth and LoRA implementations closely mirror those in their respective original publications, with optimal hyperparameters selected. Supplementary results and discussions are included in Appendix~\ref{app:settings},\ref{app:coft_oft},\ref{app:more_qualitative},\ref{app:failure}.

\textbf{Finetuning stability and convergence}. We begin by assessing the stability during finetuning as well as the convergence rate for DreamBooth, LoRA, OFT, and COFT. The outcomes of these evaluations are illustrated in Figure~\ref{fig:teaser} and Figure~\ref{fig:db_convergence}. The data indicate that both OFT and COFT maintain a high degree of stability throughout the finetuning of the diffusion model. At the 400-iteration mark, both DreamBooth and the OFT variants successfully exert control over the model output, whereas LoRA is unable to retain subject identity effectively. When extended to 2000 iterations, DreamBooth begins producing degenerate outputs, and LoRA misattributes the subject’s clothing color, depicting yellow fur instead of a yellow shirt. In contrast, OFT and COFT consistently retain stable and precise control over the generated content at this stage. These findings substantiate the efficacy of our OFT approach in facilitating both rapid and reliable convergence for subject-driven tasks. Importantly, the robustness of OFT and COFT with respect to the number of finetuning steps simplifies their practical deployment, as it reduces the need for fine-grained adjustment of the iteration count for diverse subjects. Both methods can be configured with a generous upper bound on iteration numbers without extensive hyperparameter tuning. For COFT, with an appropriately chosen $\epsilon$, setting both the learning rate and total iterations is straightforward and does not require careful tuning.

\textbf{Quantitative comparison}. In line with \cite{ruiz2023dreambooth}, we carry out quantitative evaluations to measure subject fidelity (using DINO~\cite{caron2021emerging} and CLIP-I~\cite{radford2021learning}), text prompt fidelity (using CLIP-T~\cite{radford2021learning}), and sample diversity (using LPIPS~\cite{zhang2018unreasonable}). Specifically, CLIP-I measures the mean pairwise cosine similarity between CLIP embeddings of generated images and those of the actual images. DINO operates analogously to CLIP-I, but leverages ViT S/16 DINO embeddings instead. The CLIP-T metric quantifies the mean cosine similarity between text prompt and generated image CLIP embeddings. Furthermore, we report the mean LPIPS cosine similarity between generated images for a given subject under identical text prompts. According to Table~\ref{table:dreambooth_final}, both COFT and OFT significantly surpass DreamBooth and LoRA in the DINO and CLIP-I metrics, while also matching or slightly exceeding their counterparts in prompt adherence and diversity metrics. To mitigate random variation, we repeat finetuning 30 times for each method with unique random seeds on the same pretrained model. Collectively, these results demonstrate that our OFT-based framework not only provides more stable and faster convergence, but also delivers superior and more reliable end performance.

\textbf{Qualitative comparison}. To better visualize the advantages of OFT, Figure~\ref{fig:dreambooth_final} presents a selection of random cases from our subject-driven generation experiments. For consistency across methods, each result is generated by applying the respective technique to the same finetuned model, and we avoid any tailored text prompt optimization for individual outcomes. For every approach, we report results from the version that achieves the top validation CLIP score. Reviewing Figure~\ref{fig:dreambooth_final}, it is evident that OFT and COFT are both highly effective at retaining the semantic features of the subject, whereas LoRA frequently struggles with subject identity maintenance (\emph{e.g.}, in the bowl example, LoRA fails to preserve the correct subject). Furthermore, OFT and COFT exhibit notably precise responsiveness to text guidance, whereas DreamBooth, despite successfully conserving subject identity, often generates outputs that do not align well with the textual prompt (\emph{e.g.}, see the initial row of the bowl example). These qualitative findings underscore the superiority of our OFT approach in balancing both prompt adherence and subject fidelity. Additionally, OFT demonstrates robustness to the choice of iteration count, maintaining strong performance even with extensive iterations, unlike DreamBooth or LoRA. Further visualizations are provided in Appendix~\ref{app:more_qualitative}. We also present a human study in Appendix~\ref{app:human}, which supplies further evidence of OFT’s enhanced effectiveness.

\subsection{Controllable Generation}

\textbf{Settings}. The comparative baselines used include ControlNet~\cite{zhang2023adding}, T2I-Adapter~\cite{mou2023t2i}, and LoRA~\cite{hulora2022}. We evaluate three demanding controllable generation scenarios in the main paper: canny edge to image~(C2I) conversion on the COCO dataset~\cite{lin2014microsoft}, translating segmentation maps to images~(S2I) on the ADE20K dataset~\cite{zhou2017scene}, and converting facial landmarks to portraits~(L2F) on the CelebA-HQ dataset~\cite{karras2018progressive,xia2021tedigan}. For each task, all methods fine-tune Stable Diffusion~(SD) v1.5 on the respective datasets for 20 epochs. Additional examples and analyses appear in Appendix~\ref{app:more_qualitative},\ref{app:more_control_task},\ref{app:failure}.

\textbf{Convergence}. We assess the convergence rates of ControlNet, T2I-Adapter, LoRA, and COFT on the L2F benchmark, conducting both quantitative and qualitative analyses. For quantitative assessment, we calculate the average $\ell_2$ distance between the target face landmarks and the model’s predicted face landmarks. Figure~\ref{fig:face_convergence} displays the landmark errors across different epochs of model fine-tuning. The results indicate that COFT and OFT demonstrate notably accelerated convergence relative to the other approaches. For example, LoRA requires 20 epochs to reach the level of performance that OFT achieves as early as the 8-th epoch. Importantly, both OFT and COFT employ a comparable number of trainable parameters to LoRA—significantly fewer than ControlNet—yet achieve faster convergence than these baselines. This rapid convergence by OFT is further substantiated by findings in Figure~\ref{fig:teaser}, where, notably, the rightmost sample shows that OFT reaches satisfactory convergence using only 5\% of the ADE20K training data, highlighting its superior data efficiency over both ControlNet and LoRA. In our qualitative analysis, we primarily contrast OFT, COFT, and ControlNet, given that ControlNet most closely matches our landmark error. As depicted in Figure~\ref{fig:face_convergence_qual}, both OFT and COFT consistently and smoothly align generated face poses to the prescribed control landmarks as training progresses. In contrast, ControlNet displays abrupt emergence of controllability post the 8-th epoch, a phenomenon paralleling the sudden convergence effect detailed in \cite{zhang2023adding}. This pronounced stability in the convergence behavior of OFT greatly enhances its practical applicability, as the learning process is more gradual and predictable, which in turn facilitates more reliable hyperparameter tuning for OFT.

\textbf{Quantitative comparison}. To assess controllable generation, we propose a control consistency metric. This metric involves extracting the control signal from the output image and directly comparing it with the initial input control signal. Specifically, for the C2I task, IoU and F1 score are calculated. In the case of the S2I task, we report both mean IoU and mean as well as overall accuracy. For L2F, the evaluation uses the mean $\ell_2$ distance between the ground truth control landmarks and those predicted by the model. Further information on the computation of these metrics can be found in Appendix~\ref{app:settings}. Each baseline utilizes its optimal hyperparameter configuration for evaluation. As illustrated in Table~\ref{table:control_gen}, OFT and COFT consistently deliver superior and more precise control compared to the alternative techniques. Notably, adapter-based methods (\emph{e.g.}, T2I-Adapter and ControlNet) not only converge more slowly but also achieve less favorable outcomes. LoRA exhibits stronger results on the S2I task than ControlNet, yet it lags behind in both C2I and L2F tasks. We note that the maximal attainable performance for LoRA remains relatively limited, even under well-tuned hyperparameter settings. In contrast, our empirical findings reveal that OFT continues to improve with increases in trainable parameters and its performance does not plateau quickly. Our proposed quantitative assessment for controllable generation represents an important and novel contribution, offering precise measurement of the control capabilities of finetuned models, constrained primarily by the accuracy of the segmentation or detection models employed.

\textbf{Qualitative comparison}. In addition to quantitative assessments, we qualitatively assess OFT and COFT in relation to prevailing state-of-the-art approaches, such as ControlNet, T2I-Adapter, and LoRA. The randomly sampled outputs shown in Figure~\ref{fig:control_final} illustrate that both OFT and COFT consistently deliver not only images of high realism and visual fidelity but also strong adherence to control signals. For the S2I task, it is evident that LoRA fails to respect the constraints imposed by the input segmentation map, whereas ControlNet, OFT, and COFT are all effective in maintaining control over the generative process. OFT and COFT, in particular, surpass ControlNet by synthesizing images with richer detail and more coherent geometric consistency while utilizing significantly fewer parameters. Regarding the C2I task, OFT and COFT successfully generate semantically faithful images from sparse Canny edge inputs, outperforming both T2I-Adapter and LoRA, whose results are notably inferior. For the L2F task, our approach demonstrates superior precision in pose control of generated facial images, maintaining robustness even with complex face orientations. Across all evaluated control scenarios, OFT and COFT yield images of higher qualitative quality compared to other leading baselines, attesting to the effectiveness of our OFT framework for controllable image generation. Additional qualitative results for each of the three primary control tasks can be found in Appendix~\ref{app:control_exp}. Furthermore, we show that OFT generalizes well to further control tasks—including dense pose to human body, sketch to image, and depth to image—in Appendix~\ref{app:more_control_task}.

\section{Concluding Remarks and Open Problems}

Inspired by the insight that angular relationships among neurons are pivotal to encoding visual semantics, we have introduced a straightforward yet powerful finetuning approach—orthogonal finetuning—for modulating text-to-image diffusion models. Our work focuses on two specific applications within text-to-image generation: subject-driven synthesis and controlled generation. In comparison to prevailing techniques, OFT achieves enhanced controllability and more robust finetuning stability, all while requiring a reduced number of finetuning parameters. Notably, OFT does not add to the computational cost during inference, making it well-suited for practical deployment.

Several notable open challenges are brought to light by OFT. Firstly, the method ensures orthogonality using Cayley parametrization, which necessitates computing a matrix inverse—somewhat restricting OFT’s scalability. Although adopting a block diagonal parametrization alleviates this issue to some extent, identifying methods to expedite the matrix inversion in a differentiable manner is still an open area of investigation. Secondly, OFT holds promise for compositional applications because the orthogonal matrices generated from different OFT finetuning processes can be multiplied together to yield another orthogonal matrix. Investigating whether this collection of orthogonal matrices retains the expertise for multiple downstream tasks remains an intriguing avenue for future research. Lastly, OFT’s parameter efficiency currently hinges on the use of a block diagonal design, which unavoidably brings additional bias and may constrain adaptability. Enhancing parameter efficiency through less biased and more effective architectures continues to be an important unresolved issue.